\section{Navier--Stokes: cascade, integral, and a dyadic model}\label{sec:ns}

Navier--Stokes enters this article for one structural reason, and we state the disclaimer first
and loudly: \emph{there is no connection to prime numbers here, and none is claimed}. The branch
neither feeds the twins nor feeds off them. What it shares with them is a single motto --- the
defect does not vanish, and if it does not close, it must be paid --- realised here in the most
literal currency, the viscous dissipation of energy. On this genuine analytic object we exhibit
how far the budget discipline honestly reaches, and where the millennium problem stands untouched.
The equation is formalised in the strong form on $E_3 := \code{EuclideanSpace}\ \R\ (\code{Fin }3)$,
with the differential operators (divergence, vector Laplacian, convective term) assembled from bare
\code{fderiv}.

\begin{definition}[\code{IsNSSolution}]\label{def:IsNSSolution}
For a viscosity $\nu$, an external force $f$, a velocity field $u$ and a pressure $p$,
\begin{equation}\label{eq:ns-solution}
\code{IsNSSolution}\ \nu\ f\ u\ p \;:=\;
\bigl[\partial_t u + (u\cdot\nabla)u = \nu\,\Delta u - \nabla p + f\bigr]
\;\wedge\; \bigl[\operatorname{div} u = 0\bigr].
\end{equation}
\end{definition}

\noindent These are the incompressible Navier--Stokes equations \cite{clay-ns,millennium} in the
classical strong form. No weak solutions and no distributions: the strong form is chosen
deliberately, so that the predicate can be read with the eyes. The first question to any such
predicate is whether it is empty; the answer is by machine.

\begin{theorem}[\code{zero\_is\_NSSolution}]\label{thm:zero_is_NSSolution}
$\code{IsNSSolution}\ \nu\ 0\ 0\ 0$ for every $\nu \in \R$: the zero field with zero pressure and
no force is a solution at any viscosity. \green{}
\end{theorem}

\emph{Proof idea.} Direct computation; every term of \cref{eq:ns-solution} vanishes. Modest, but
this is precisely the inoculation against vacuity: the predicate is inhabited, the equation is real.

The energetics is given by Bochner integrals over the volume: the kinetic energy
$\code{kineticEnergy} = \tfrac12\int\|u\|^2$ and the dissipation rate
$\code{dissipationRate} = \nu\int\sum_i\|\partial_i u\|^2$ (the enstrophy form). A trap hides here,
and the repository exposes it by machine: the Bochner integral in \code{mathlib} is
\emph{silently equal to zero} on a non-integrable function (\code{MeasureTheory.integral\_undef}).

\begin{theorem}[\code{kineticEnergy\_of\_not\_integrable}]\label{thm:kineticEnergy_of_not_integrable}
If $\|u\|^2$ is not integrable (the predicate \code{FiniteKineticEnergy} fails), then
$\code{kineticEnergy}\ u = 0$. \green{}
\end{theorem}

\emph{Proof idea.} A rewrite through \code{MeasureTheory.integral\_undef}. The equality
$\code{kineticEnergy}\ u = 0$ may therefore mean not ``the energy is zero'' but ``the energy is
infinite or undefined''; the ``zero energy'' of a non-integrable field is an artefact of the
definition, not physics. The moral is that every energy statement is obliged to travel in a pair
with a named integrability predicate --- here this is not a wish but a proven warning.

\paragraph{From a monolith to a pointwise balance.}
The analytic core of the branch is the energy inequality
$E(u(t_2)) + \int_{t_1}^{t_2} D(u(s))\,ds \le E(u(t_1))$, named \code{TwoTimeEnergyInequality}. The
work did not prove it; it \emph{narrowed} it to a single pointwise input \code{EnergyBalanceLaw},
the \code{HasDerivAt} identity $dE/dt = -D(t)$ --- the classical energy balance of a smooth rapidly
decaying solution, in which convection and pressure do no work thanks to $\operatorname{div} u = 0$.
The glue from the narrow input is proved.

\begin{theorem}[\code{twoTimeEnergyInequality\_of\_energyBalance}]\label{thm:twoTimeEnergyInequality_of_energyBalance}
If the pointwise balance $\code{EnergyBalanceLaw}\ \nu\ u$ holds and $s \mapsto D(u(s))$ is
interval-integrable on every $[t_1,t_2]$, then $\code{TwoTimeEnergyInequality}\ \nu\ u$ holds. \green{}
\end{theorem}

\emph{Proof idea.} The pointwise balance plus integrability yield the two-time inequality (in fact
an equality) via the fundamental theorem of calculus
(\code{intervalIntegral.integral\_eq\_sub\_of\_hasDerivAt}). The input has become strictly narrower:
not an integral inequality over all pairs of times, but a single differential identity.

\paragraph{The chain up to the cascade.}
Let $\code{DissipativeStage}\ \nu\ u\ \delta\ t_1\ t_2$ mark a time interval on which at least
$\delta$ of dissipation has accumulated. The budget machine forbids an infinite tower of such steps.

\begin{theorem}[\code{ns\_no\_infinite\_dissipative\_cascade}]\label{thm:ns_no_infinite_dissipative_cascade}
Let $\delta > 0$ and suppose $\code{TwoTimeEnergyInequality}\ \nu\ u$ holds. Then there is no
sequence of times $(t_k)_{k\in\N}$ with $\code{DissipativeStage}\ \nu\ u\ \delta\ t_k\ t_{k+1}$ for
all $k$. \green{} (conditional on the inequality input)
\end{theorem}

\emph{Proof idea.} The accumulated payment $\sum_k \delta$ would exceed the starting energy
$E(u(t_0))$; a direct application of the quantised cascade lemma
\code{no\_infinite\_uniform\_dissipative\_cascade}. This is exactly the certificate that regularity
demands --- \emph{if} one is handed the premise. The word ``quantised'' carries all the weight, and
the reason is recorded by machine.

\begin{theorem}[\code{real\_positive\_work\_not\_wellfounded}]\label{thm:real_positive_work_not_wellfounded}
There exists $a : \N \to \R$ with $a_{n+1} < a_n$, $a_n > 0$ and $a_n - a_{n+1} > 0$ for all $n$:
an $\R$-sequence with strictly positive work at every step yet a finite total descent. \green{}
\end{theorem}

\emph{Proof idea.} The banal witness $a_n = 1/2^n$. For the $\N$-engine of the twins strict positive
work forces well-foundedness for free; for $\R$ it does not. Hence the analogy with the twin engine
is only structural: the finiteness of the cascade is bought by a uniform lower bound $\delta > 0$,
not by positivity as such. Where to obtain $\delta$ for a real solution is analysis, not structure,
and it stays open.

\paragraph{The integral, taken.}
The FTC glue lifts to an exact identity, and the solution itself is recovered by integration.

\begin{theorem}[\code{energy\_identity\_of\_energyBalance}]\label{thm:energy_identity_of_energyBalance}
Under $\code{EnergyBalanceLaw}\ \nu\ u$ and interval-integrability of the dissipation, for all
$t_1,t_2$ one has
\begin{equation}\label{eq:ns-energy-identity}
E(u(t_2)) = E(u(t_1)) - \int_{t_1}^{t_2} D(u(s))\,ds,
\end{equation}
as an equality; the two-time inequality of
\cref{thm:twoTimeEnergyInequality_of_energyBalance} is a coarsening of it. \green{}
\end{theorem}

A hypothetical finite-time singularity is a \code{SingularCascade}: infinitely many
$\delta$-dissipative stages compressed towards a single finite moment $T$ --- the same perpetual
engine as elsewhere. Its impossibility is green.

\begin{theorem}[\code{noSingularCascade\_of\_energyBalance}]\label{thm:noSingularCascade_of_energyBalance}
The pointwise energy balance $dE/dt = -D$ together with integrability implies that no rupture points
of cascade type exist: $\code{NoSingularCascade}\ \nu\ u$. \green{}
\end{theorem}

\emph{Proof idea.} A direct application of \cref{thm:ns_no_infinite_dissipative_cascade}: the budget
forbids an infinite $\delta$-cascade, and one compressed to a finite time all the more. We stress
the honest hedge: \code{NoSingularCascade} is \emph{not} $C^\infty$-regularity but the absence of an
infinite uniformly quantised dissipative ladder --- a surrogate. That the qualifier ``uniformly''
is load-bearing is shown by a forged non-uniform cascade
(\code{cookedProfileCascade\_not\_uniform}, \green{}), whose drops slip under every $\delta$ and
which the budget cannot take.

\paragraph{Down to $\R^3$: the integral assembled on the box.}
Above, the energy lived as an abstract profile handed over by hypothesis. On the class of
box-supported solutions the balance is no longer postulated but \emph{derived}, engaging
\code{mathlib}'s divergence theorem (in its box form
\code{integral\_divergence\_of\_hasFDerivAt\_off\_countable}) and differentiation under the integral
sign, through three integration-by-parts identities (pressure and convection die by
divergence-freeness; viscosity yields minus the enstrophy).

\begin{theorem}[\code{noSingularCascade\_of\_r3Assembly}]\label{thm:noSingularCascade_of_r3Assembly}
For a box-supported $C^2$ solution with force $f = 0$, cascade smoothness $\code{NoSingularCascade}$
is derived in green, without the first-cause decree \code{nsBoundary}. \green{}
\end{theorem}

\emph{Proof idea.} On the box, \code{energyBalance\_of\_boxSupported} derives $dE/dt = -D$ from the
derivative under the integral and the three integration-by-parts identities, after which the budget
machine of \cref{thm:noSingularCascade_of_energyBalance} takes over; \code{\#print axioms} shows the
standard triple, no \code{sorry}, no new axiom. Beyond the box class the limiting passage
box $\to \R^3$ for merely decaying (non-compactly-supported) fields is absent from \code{mathlib} ---
an external gap --- and \code{EnergyBalanceLaw} remains an open \red{} input in general.

\paragraph{The Clay reduction: one missing criterion.}
The exact Clay statement (A) --- Schwartz data, divergence-freeness, the existence of a smooth
global finite-energy solution --- is encoded and machine-reduced to a single open statement.

\begin{theorem}[\code{clayA\_of\_regularityTransfer\_and\_vorticityControl}]\label{thm:clayA_of_regularityTransfer_and_vorticityControl}
The known transfer theory --- Kato's local strong existence \cite{kato1984} together with the
Beale--Kato--Majda continuation criterion \cite{bkm1984} --- and the single criterion
\code{GlobalVorticityControl} imply Clay-(A). \green{}
\end{theorem}

\emph{Proof idea.} Pure modus ponens: a local smooth solution exists; the criterion keeps the
vorticity bounded on each finite interval; \cite{bkm1984} extends the solution to all time. The
theorem is green and \emph{not vacuous} --- the conclusion is inhabited (\code{globalSmoothSolution\_zero})
and both inputs are genuinely consumed. The point is isolation: everything left for (A) is pulled
into one named open core, \code{GlobalVorticityControl}, the supercriticality barrier that no one
knows how to prove. The transfer theory itself (\code{RegularityTransfer}, \code{RegularityNecessity})
is honestly named as a \red{} input --- real theorems of the literature \cite{kato1984,bkm1984} not
yet formalised in \code{mathlib}. It is tempting to close (A) through the engine, as one can for
Collatz; for Navier--Stokes that bridge provably tears, and the module records this by machine.

\begin{theorem}[\code{greenBudget\_strictly\_weaker\_than\_vorticityControl}]\label{thm:greenBudget_strictly_weaker_than_vorticityControl}
The green budget machine is strictly weaker than \code{GlobalVorticityControl}: it does not imply
the criterion. \green{}
\end{theorem}

\emph{Proof idea.} The surrogate \code{NoSingularCascade} kills only the uniformly quantised
cascade, while \code{cookedProfileCascade\_not\_uniform} exhibits a non-uniform cascade that slips
away. A genuine blow-up is a continuum phenomenon; the discrete engine prohibition does not reach
it. Honesty demands naming the open core analytically (the vorticity), not hiding it in a metaphor.
This work therefore brings the Navier--Stokes solution \emph{not one step closer} mathematically:
it gives the exact formulation, a verified reduction skeleton, the exact name of the single missing
theorem, and the barrier that keeps it out of reach. Navier--Stokes is \emph{not} solved.

\paragraph{A dyadic model: where the engine works, and where it blows up.}
To test the engine reading against a rigorous discrete fluid, we formalise the Katz--Pavlovi\'c
dyadic model \cite{katz-pavlovic2005}. (No fluid model had been formalised in any proof assistant
before; the novelty is formalisational, not mathematical --- nothing here is solved.) Where the
dissipation is uniform, the budget wins; where the cascade is genuine, the engine is
\emph{realised}, and the naive ``no engine $\Rightarrow$ no singularity'' is refuted. The rigorous
core is an ODE fact.

\begin{theorem}[\code{superlinear\_blowup\_sq}]\label{thm:superlinear_blowup_sq}
No global positive $C^1$ function can satisfy $y'(t) \ge C\,y(t)^2$ with $C > 0$: the assumption of
a global solution yields \code{False}. \green{}
\end{theorem}

\begin{proof}
Put $w(t) := 1/y(t) + C\,t$. Then $w'(t) = -y'/y^2 + C \le -C + C = 0$, so $w$ is non-increasing,
whence $1/y = w - C\,t$ must go negative in finite time, contradicting $y > 0$. A super-linear
cascade is a realised perpetual engine: it exists exactly until the moment of blow-up.
\end{proof}

\begin{theorem}[\code{dyadic\_blowup}]\label{thm:dyadic_blowup}
The Katz--Pavlovi\'c model \code{DyadicSolution}, with equations
$a_n' = \lambda^n a_{n-1}^2 - \lambda^{n+1} a_n a_{n+1} - d_n a_n$, is globally empty: blow-up is
inevitable. \green{}
\end{theorem}

\emph{Proof idea.} The nonlinear transfer telescopes (energy is conserved by the coupling), and the
blow-up follows from \cref{thm:superlinear_blowup_sq} applied to the leading mode. The super-linear
drive $y' \ge C\,y^2$ is no longer postulated: for a whole class of solutions it is derived from the
$\lambda^n$-couplings.

\begin{theorem}[\code{frontDrive\_of\_invariant}]\label{thm:frontDrive_of_invariant}
Suppose a Katz--Pavlovi\'c solution has one front shell pinched from below by its two neighbours ---
the invariant \code{FrontDomination}: $\rho\, a_{J+1} \le a_J$, $a_{J+2} \le \kappa\, a_{J+1}$,
$m \le a_{J+1}$. Then the drive $C\,y^2 \le y'$ is derived directly from the $\lambda^n$-couplings,
with $C = \lambda^{J+1}\rho^2 - \lambda^{J+2}\kappa - d_{J+1}/m$. \green{}
\end{theorem}

\emph{Proof idea.} The three estimates are substituted into the right-hand side of shell $J+1$: the
inflow $\lambda^{J+1} a_J^2$ produces $y^2$, and the outflow and dissipation are subtracted under
control. Non-vacuity is confirmed by the exact self-similar profile
$a_n(t) = \lambda^{-n}/((\lambda^2-1)(T-t))$, which satisfies the invariant. One input remains open:
the preservation of \code{FrontDomination} for \emph{infinitely} many modes over the whole lifespan
--- the moving Katz--Pavlovi\'c front, a research frontier, not proved here.

The origin of the cascade carries the only deliberate yellow layer of this branch. The bottom shell
$n = 0$ has zero inflow: it can only give, never start itself, so its origin cannot be begotten from
inside the couplings. The self-similar solution needs an external pumping term, and that supply is
drawn from the first cause.

\begin{theorem}[\code{dyadicBlowup\_is\_firstCauseManifestation}]\label{thm:dyadicBlowup_is_firstCauseManifestation}
The supply at scale $n = 0$ is drawn from the boundary \code{nsBoundary} of the axiom \axiom{}: the
dyadic blow-up joins the first-cause manifestations through its origin. \yellow{}~(\textsc{axiom-tainted})
\end{theorem}

\emph{Proof idea.} Two walls: a green wall of impossibility (the origin is uncausable from inside)
plus a supply decreed from outside. The green derivation of the drive is self-contained and does not
depend on this layer --- remove the first cause and the mathematics of the blow-up stays intact,
only the answer to ``who lit the pump $n=0$'' disappears --- which is why the taint grows by exactly
two declarations and not one more.

\paragraph{Honest status.} The engine reading is neither vindicated nor discarded: it is
\emph{mapped}. Where dissipation is uniform the budget forbids a perpetual cascade; where the
cascade is genuine the engine is realised and the fluid blows up in the model. No new mathematics
for the millennium problem is produced (this is the shadow of Tao's supercriticality barrier
\cite{clay-ns}), only the world's first formalisation of a discrete fluid model of \emph{known}
mathematics, an exact reduction skeleton, and a machine map of where the reading is valid and where
it breaks. Navier--Stokes is \emph{not} solved.
